%%
%% Front matter: rights, conference, title, authors, abstract, keywords.
%% Included by main.tex before \maketitle.

%%
%% Rights management information.
\copyrightyear{2026}
\copyrightclause{Copyright for this paper by its authors.
  Use permitted under Creative Commons License Attribution 4.0
  International (CC BY 4.0).}

%%
%% Conference information.
\conference{CLEF 2026 -- Conference and Labs of the Evaluation Forum,
  September 21--24, 2026, Jena, Germany}

%%
%% Title.
\title{Fine-Tuning of BioCLIP 2.5 with Taxonomic Heads for Multi-Species Plant Identification}

%%
%% Authors.
%% Per the canonical CEUR-ART template, each author block supports an
%% optional key-value list with \texttt{orcid=}, \texttt{email=},
%% \texttt{url=}. Only \texttt{email=} is supplied here; add the other
%% keys when convenient.
\author[1]{Arjun Raj}[%
  email=u7526852@anu.edu.au,
]
\fnmark[1]

\address[1]{The Australian National University, Canberra, ACT 2601, Australia}

\author[1]{Manindra de Mel}[%
  email=u7543337@anu.edu.au,
]
\cormark[1]
\fnmark[1]

\author[1]{Razeen Wasif}[%
  email=u7619822@anu.edu.au,
]
\fnmark[1]

\author[1]{William Brake}[%
  email=u7480294@anu.edu.au,
]
\fnmark[1]

\cortext[1]{Corresponding author.}
\fntext[1]{These authors contributed equally.}

%%
%% Abstract.
\begin{abstract}
We describe ANU's submission to PlantCLEF 2026, where the task is to
predict every species present in a vegetation-quadrat image. The
central challenge is a domain shift where training images show single
plants while test images show dense multi-species plots; this is
further compounded by a long-tailed label distribution. Our system is
BioCLIP~2.5 ViT-H/14 where only the last four transformer blocks plus
the final layer norm and projection are unfrozen, while the lower
twenty-eight blocks stay frozen to preserve the Tree-of-Life prior. We
train on an enlarged manifest that combines PlantCLEF~2024 with a
research-grade iNaturalist pull (2.65~million single-plant images,
7{,}806 species) with auxiliary genus- and family-level cross-entropy
heads that regularise the species head without adding inference cost. At inference each quadrat is split into a $4{\times}4$ grid of
sixteen tiles, each tile is forwarded through the encoder at both
$224$ and $336$ pixels, the per-tile softmax distributions are
averaged across tiles and across resolutions, class-prior logit
adjustment is applied, and an adaptive probability threshold selects
the emitted species set. Our best public Macro F1 is \textbf{0.41826}; the same
configuration scores \textbf{0.40283} private (the strongest private
result of the project, $0.40600$, comes from a logit-adjusted
single-resolution variant). Project website: \url{https://arm-wision.github.io/}.
\end{abstract}

%%
%% Keywords.
\begin{keywords}
  plant identification \sep
  multi-label classification \sep
  partial fine-tuning \sep
  BioCLIP \sep
  taxonomic auxiliary loss \sep
  vegetation plot analysis \sep
  PlantCLEF 2026
\end{keywords}
